\section{Технологическая часть}
\subsection{Средства реализации}

Для реализации программного обеспечения выбран язык программирования C++20~\cite{cpp-standard}.
Выбор языка обусловлен:
\begin{itemize}
    \item наличием опыта работы с языком;
    \item наличием стандартной библиотеки (STL), которая предоставляет контейнеры для работы с данными;
    \item возможностью использования нативных потоков;
    \item поддержкой объектно-ориентированной парадигмы;
    \item кроссплатформенностью.
\end{itemize}

Для реализации многопоточной обработки применяются средства стандартной библиотеки C++~\cite{cpp-standard}: \texttt{std::thread} для создания потоков, \texttt{std::mutex} и \texttt{std::condition\_variable} для синхронизации, \texttt{std::future} и \texttt{std::packaged\_task} для асинхронного программирования, контейнеры \texttt{std::vector}, \texttt{std::array} и умные указатели \texttt{std::shared\_ptr} для управления памятью.

Графический интерфейс реализован с использованием фреймворка Qt 6~\cite{qt6-doc}.
Применяются модули: \texttt{QtWidgets} для создания оконного интерфейса, \texttt{QtConcurrent}
для параллельного выполнения задач, \texttt{QFutureWatcher}
для отслеживания асинхронных операций, \texttt{QPixmap} и \texttt{QPainter} для работы с изображениями.
Основные компоненты интерфейса: \texttt{QMainWindow}, \texttt{QLabel}, \texttt{QPushButton}, \texttt{QFileDialog}.
Qt обеспечивает безопасную многопоточность через механизм сигналов и слотов.

Разработка выполнена в среде разработки Visual Studio Code~\cite{vscode}. Для сборки использовался CMake~\cite{cmake}.
Для модульного тестирования применяется фреймворк Google Test~\cite{gtest}.

\subsection{Реализация алгоритмов}
Реализация функции алгоритма трассировки пути представлены на листинге~\ref{b-listing_1}

\clearpage
